\chapter{建模工具進階}\label{ch:modeling-tools}\label{sec:modeling-reference}

前兩章透過具體模型介紹常見的建模過程。本章按操作目標說明各項工具，便於在編輯時查閱引數和檢查方法。選擇、晶胞處理、軟物質構建和過程記錄各有相應入口，可以從正在處理的問題開始閱讀。

建模頁面隨當前文件改變。分子中的變換和測量主要在\ui{分子建模}的\ui{原子編輯}、\ui{幾何編輯}中；晶體中的區域性分子使用\ui{晶體建模}的\ui{組分變換}、\ui{組分幾何}。改變晶體整體時則使用\ui{晶格編輯}或\ui{超胞與晶格}。命令呈灰色時，先檢查當前文件型別和選擇數量。例如，\ui{移動原子}要求一個獨立位點，\ui{設定鍵角}要求依次指定三個原子。

\section{精確編輯結構}

\subsection{儲存並再次呼叫一組原子}\label{sec:selection-set-workflow}
\index{Xuanzeji@選擇集}

選擇集為一組模型位點命名，適合反覆處理表面底層、吸附分子或某一段鏈。呼叫選擇集後，可再次選中這些位點。

\begin{steps}
  \item 在三維檢視中選中目標。單擊選一個原子，Shift 或 Ctrl 配合單擊可調整多選；\ui{B}進入框選，\ui{L}進入套索。若要選擇整個分子，可雙擊一個原子；也可在原子右鍵選單中選擇\ui{Select Connected Component}。連通關係取決於當前顯示模型使用的鍵，晶體網路可能因此選中很大一片結構。
  \item 若要選擇一種元素，在該元素的原子上開啟右鍵選單，選擇\ui{Select All \ldots\ Atoms}。先檢視選擇數，再開始替換或刪除。
  \item 開啟頂部建模頁面的\ui{選擇集}，單擊\ui{儲存當前選擇}，輸入能說明用途的名稱，例如“底層固定原子”。
  \item 單擊空白區域或按 Esc 清除選擇，再從\ui{選擇集}選單中單擊剛儲存的名稱，確認相同位點重新亮起。
  \item 儲存專案。若還要保留這組選擇的完整模型資訊，可同時匯出原生結構 JSON；CIF、XYZ 等交換格式並不承擔儲存全部選擇資訊的用途。
\end{steps}

選單名稱後的兩個數表示當前還能找到的位點數與原先儲存的位點數。例如，某組中的原子被刪除後，前一個數會變小；召回時也會提示有位點已不存在。不能把剩餘選擇當成最初的完整一組。對同名選擇重新執行儲存，會更新該組內容；選單提供選擇集的儲存與召回。

週期結構應額外核對\ui{sites}。顯示重複範圍改變後，同一位點可能出現多個映像；選擇集仍按模型位點工作，不會因為螢幕上的球變多而自動增加成員。

\subsection{修正原子的種類與位置}\label{sec:atom-edit-reference}

\ui{新增原子}直接加入一個位點。晶體可以選擇分數座標或笛卡爾座標，分子使用以\AA 為單位的笛卡爾座標。與之不同，分子中的\ui{繪製原子}還可以同時指定連線物件與鍵級；要建立化學連線，應使用草繪或新增鍵，而不是隻把兩個原子放得很近。

\begin{steps}
  \item 在\ui{原子編輯}中開啟\ui{新增原子}，輸入元素，確認座標型別，再分別填寫三個座標分量。先預覽，確認沒有放到已有原子上，再應用。
  \item 要替換已有原子，先選擇它們，再開啟\ui{替換元素}，填寫目標元素。替換會保留這些位點的位置，不會重新搜尋適合該元素的幾何結構。
  \item 要準確移動一個位點，開啟\ui{移動原子}，輸入該原子的目標位置。晶體中先切換座標型別再輸入，避免把分數座標當成\AA。
  \item 要讓一組原子一起移動，使用\ui{平移原子}。例如，笛卡爾向量 \((0,0,1)\) 表示整體沿全域性 $z$ 軸移動 $1\,\text{\AA}$；晶體中的分數向量 \((0,0,1)\) 表示移動一個 $\mathbf c$ 晶格向量。
  \item 刪除時先核對獨立位點數，開啟\ui{刪除原子}並檢查序號。應用後檢查組成和剩餘連線，最後另存結構或專案。
\end{steps}

這些區域性晶體命令還可能出現\ui{保留對稱性}或\ui{打破對稱性}的選擇。保持對稱性會按等價關係處理相關位點；做區域性摻雜或缺陷時，應選擇與實際目的相符的方式。執行前看預覽中的所有變化，不能只觀察最初點中的那一個原子。

\subsection{整理重複位點和原子順序}\label{sec:atom-clean-reference}

\ui{合併原子}處理相距非常近的晶體位點。它按以\AA 為單位的容差比較週期距離，不是把重複顯示的邊界映像變成更少的球。預設容差為 $0.01\,\text{\AA}$；初次處理匯入模型時，應先用較小的容差檢查。

\begin{steps}
  \item 儲存原始模型副本，開啟\ui{原子編輯}中的\ui{合併原子}。
  \item 選擇合併方式。\ui{刪除重合原子}保留一份組分資訊；\ui{佔位率相加}會合並同簇位點的物種佔位；\ui{屬性取平均}用於可相容的數值位點屬性。各方式都會把同簇位置移到按鄰近週期映像求得的中心，不能理解為座標一定保留不變。
  \item 預覽後檢查位點數、佔位和區域性位置。若兩個真實相鄰原子被併成一個，取消並減小容差。
  \item 應用後儲存為新名稱，保留原輸入供比較。
\end{steps}

\ui{原子排序}按元素重新排列位點，\ui{倒序排列}改變排序方向。排序用於整理檔案和序號，不改變原子的空間位置。已經記下的“第幾個原子”應重新核對；需要長期追蹤的一組位點更適合用選擇集儲存。

\ui{擾動原子}和\ui{擾動結構}提供隨機幾何擾動，引數是最小、最大位移，單位均為\AA，以及隨機種子。例如，可用最小位移 0、最大位移 $0.05\,\text{\AA}$、種子 7 做一次很小的擾動練習。這兩個命令作用於整個模型，包括未選中的原子。預覽時檢視整體變化，並保留原始副本供比較。

\subsection{改變分子的方向}\label{sec:orientation-reference}

居中、旋轉和映象用於改變幾何擺放。它們本身不執行能量最佳化。晶體中的吸附物也可使用\ui{組分變換}進行這些操作，但應先選擇完整組分。

\subsubsection{把一組原子移到目標中心}

居中時，同一組中的所有原子使用同一個位移。設選擇中有 $N$ 個原子，分子用笛卡爾位置 $\mathbf r_i$，晶體用分數座標 $\mathbf f_i$，兩種情況分別為
\begin{equation}
 \begin{aligned}
 \Delta\mathbf r&=\mathbf r_{\mathrm{target}}-
        \frac1N\sum_{i=1}^{N}\mathbf r_i,
 &\mathbf r_i'&=\mathbf r_i+\Delta\mathbf r,\\
 \Delta\mathbf f&=\mathbf f_{\mathrm{target}}-
        \frac1N\sum_{i=1}^{N}\mathbf f_i,
 &\mathbf f_i'&=\mathbf f_i+\Delta\mathbf f.
 \end{aligned}
 \label{eq:center-selected-atoms}
\end{equation}
第一行用於分子，位置與位移的單位都是\AA；第二行用於晶體，各量均為無量綱分數座標。分子常用目標 $(0,0,0)$，晶體常用 $(0.5,0.5,0.5)$。每個原子在平均值中佔相同權重，因此居中不使用原子質量，也不改變所選原子之間的距離。

例如，分子中兩個原子的 $x$ 座標為1和3\,\AA，要把它們的平均位置移到原點，共同位移就是 $-2\,\text{\AA}$；處理後為 $-1$ 和1\,\AA，間距仍為2\,\AA。晶體中若兩個位點的第一分數分量為0.1和0.3，目標為0.5，則都加0.3，得到0.4和0.6。這裡只取一個分量說明計算，其餘兩個分量同樣處理。

週期邊界附近要先看清使用的是哪一個映像。0.99與0.01的普通平均為0.5，但兩個位點可能實際聚集在邊界兩側；若要以這段連續結構的中心為準，應先整理為相鄰的連續映像再居中。應用後坐標可能被歸回晶胞，因此還應從週期連線看整體，而不只比較某個座標的正負。

\subsubsection{繞指定錨點旋轉}

旋轉先從每個位置中減去錨點，再轉動這個相對向量，最後加回錨點。這樣，錨點本身不動，離錨點的距離也不變。沿用行向量表示，有
\begin{equation}
 \mathbf r'=\mathbf p+(\mathbf r-\mathbf p)
              R(\hat{\mathbf n},\theta)^{\mathsf T}.
 \label{eq:rotate-row-position}
\end{equation}
$\mathbf p$ 是以\AA 為單位的錨點，$\hat{\mathbf n}$ 是旋轉軸的單位方向，$\theta$ 是轉角；上標 $\mathsf T$ 表示轉置。對話方塊中填寫度數，正角遵循右手定則：右手拇指沿輸入軸的正向，四指彎曲的方向為正轉向。軸向量只決定方向，例如\literalinput{0 0 1}和\literalinput{0 0 2}表示同一根正向轉軸。

繞全域性 $z$ 軸旋轉時，常用的矩陣為
\[
 R_z(\theta)=\begin{bmatrix}
 \cos\theta&-\sin\theta&0\\
 \sin\theta&\cos\theta&0\\
 0&0&1
 \end{bmatrix}.
\]
取錨點 $\mathbf p=(1,1,0)\,\text{\AA}$，位置 $\mathbf r=(2,1,0)\,\text{\AA}$，繞 $z$ 軸正轉 $90^\circ$，得到 $\mathbf r'=(1,2,0)\,\text{\AA}$。這個小例子可用來辨認旋轉方向；旋轉後原子到錨點仍相距1\,\AA。晶體中原子的等價位置還會歸回晶胞，預覽時應連同邊界映像一起觀察。

\subsubsection{關於一個平面作映象}

映象需要一個平面上的點 $\mathbf p$ 和非零笛卡爾法向 $\mathbf n$。先把法向歸一化，再減去兩倍的有向投影：
\begin{equation}
 \begin{aligned}
 \hat{\mathbf n}&=\frac{\mathbf n}{\|\mathbf n\|},\\
 \mathbf r'&=\mathbf r-
 2\bigl[(\mathbf r-\mathbf p)\cdot\hat{\mathbf n}\bigr]
       \hat{\mathbf n}.
 \end{aligned}
 \label{eq:mirror-cartesian-plane}
\end{equation}
點積 $(\mathbf r-\mathbf p)\cdot\hat{\mathbf n}$ 是原子到平面的有向距離，單位為\AA。映象後它的符號反轉，大小相同；位於平面上的原子保持不動。法向填\literalinput{1 0 0}，平面點取 $(1,0,0)\,\text{\AA}$，就是 $x=1\,\text{\AA}$ 的平面。位置 $(1.5,2,0)\,\text{\AA}$ 映象後變為 $(0.5,2,0)\,\text{\AA}$，兩側到平面均為0.5\,\AA。

\begin{figure}[!htb]
\centering
\begin{tikzpicture}[x=1mm,y=1mm,font=\sffamily\small]
 \node[text=brand] at (24,47) {绕锚点旋转};
 \draw[teal,dashed] (10,15)--(33,15);
 \draw[teal,dashed] (10,15)--(10,38);
 \draw[->,>=stealth,brand,line width=.85pt] (33,15) arc[start angle=0,end angle=90,radius=23];
 \fill[ink] (10,15) circle (1.1);
 \fill[teal] (33,15) circle (1.4);
 \fill[brand] (10,38) circle (1.4);
 \node[anchor=north] at (10,13) {$\mathbf p$};
 \node[anchor=west,text=teal] at (35,15) {$\mathbf r$};
 \node[anchor=east,text=brand] at (8,38) {$\mathbf r'$};
 \node[text=brand] at (30,35) {$+90^\circ$};
 \node[text=muted,font=\sffamily\footnotesize] at (24,3) {到锚点的距离保持不变};
 \node[text=brand] at (86,47) {关于平面镜像};
 \fill[paleblue] (83,8) rectangle (87,40);
 \draw[teal,line width=.8pt] (85,8)--(85,40);
 \draw[dashed,muted] (68,29)--(102,29);
 \draw[->,>=stealth,brand,line width=.7pt] (85,12)--(105,12) node[above] {$\hat{\mathbf n}$};
 \fill[ink] (85,18) circle (1.1);
 \fill[teal] (102,29) circle (1.4);
 \fill[brand] (68,29) circle (1.4);
 \node[anchor=south,text=teal] at (102,31) {$\mathbf r$};
 \node[anchor=south,text=brand] at (68,31) {$\mathbf r'$};
 \node[anchor=east] at (83,18) {$\mathbf p$};
 \node[text=muted,font=\sffamily\footnotesize] at (86,3) {到镜面的距离两侧相等};
\end{tikzpicture}
\caption{锚点旋转与平面镜像（截面示意）。左图中的点绕 $\mathbf p$ 旋转；右图中的点跨过镜面，沿法向到镜面的距离大小不变。}
\label{fig:round10-anchor-mirror}
\end{figure}


居中、旋轉和映象均為幾何操作。處理分子時，先選擇需要整體移動的部分，再按下面的步驟操作；若只需調整觀察方向，旋轉相機即可。


\begin{steps}
  \item 選擇要處理的一組原子，開啟\ui{原子居中}。分子的目標中心以\AA 表示，常用 \((0,0,0)\)；晶體的目標中心為分數座標，常用 \((0.5,0.5,0.5)\)。這裡的中心是所選原子位置的算術平均，不是質量中心。
  \item 要作已知角度的轉動，開啟\ui{旋轉原子}。角度用度，旋轉軸為全域性笛卡爾方向，錨點用\AA。例如，軸 \((0,0,1)\)、錨點 \((0,0,0)\)、角度 $90^\circ$ 表示繞過原點的 $z$ 軸轉動。
  \item 要映象，開啟\ui{映象原子}。填寫平面法向和一個平面內點。例如，法向 \((1,0,0)\)、平面點 \((0,0,0)\) 定義 $yz$ 平面。法向只確定方向，不是一個需要輸入\AA 的長度。
  \item 預覽時比較原子間距離、連線與附近宿主。選完整分子作剛性變換時，內部分子幾何應保持不變；只選一部分則可能拉斷或扭曲原有幾何。
  \item 應用並另存。若只是調整觀看方向，退出建模工具後旋轉相機即可，無須改變結構座標。
\end{steps}

\ui{幾何對齊}適合把長分子的主軸放到指定方向，或讓一組近似共面的原子的平面法向對齊。選擇\ui{principal\_axis}或\ui{plane\_normal}，再輸入非零目標向量。前者看整體主方向，後者看擬合平面；對高度對稱的分子，主軸可能不是唯一的，完成後仍應從不同方向檢視。

\subsection{修改顯式化學鍵}\label{sec:bond-edit-reference}

顯式化學鍵記錄哪兩個原子相連及其鍵級；三維顯示中的線條還可能來自按距離推斷的鄰接關係。修改拓撲時，應以分子中的連線記錄為準。鍵級 1、2、3 分別表示單、雙、三鍵，1.5 用於這裡的芳香連線。

\begin{steps}
  \item 匯入或開啟要修改的分子，儲存一份副本。依次選擇構成目標鍵的兩個原子，檢視序號，不能多選。
  \item 兩原子尚未連線時，在\ui{分子建模}選單中開啟\ui{新增化學鍵}，指定鍵級後應用。
  \item 已有連線需要改變鍵級時，使用\ui{設定鍵級}；需要斷開連線時，使用\ui{刪除化學鍵}。刪除鍵保留兩個原子，刪除原子則會改變原子數，兩者應按目的區分。
  \item 在\ui{原子化學屬性}中核對受影響原子的元素、形式電荷、雜化與芳香標記。不替換元素時保留\ui{unchanged}。形式電荷是選定原子的屬性，不是把整個分子的淨電荷平均分配到所有原子。
  \item 若原有氫已不適合修改後的連線，選擇相應重原子後先用\ui{刪除氫}，再用\ui{新增氫}按新連線補齊。核對組成、重原子連線和氫數，並測量可疑的短距離；必要時再整理區域性幾何。
  \item 儲存專案與原生結構 JSON，保留鍵級和原子屬性。芳香分子的隨書例子採用 JSON；不要僅憑某種交換格式中顯示出了線條，就認定全部芳香資訊已經保留。
\end{steps}

\ui{分子診斷}可以執行規則檢查，結束後顯示\ui{Molecule diagnostics completed.}。圖形操作中的檢查按第~\ref{ch:modeling}~章\ref{sec:molecule-manual-check}節進行：先看組成和連線，再測量區域性幾何；MATLAB\regmark{}工具箱使用者可從診斷返回資料讀取逐項記錄。

直接草繪時，按 S 進入工具，面板可指定元素、鍵級、形式電荷和雜化提示。從已有原子拖出可形成帶鍵原子，拖到另一個已有原子可嘗試建立鍵；空白處單擊則放置原子。向已有分子新增連線時，先選中宿主連線位點。出現碰撞提示時應移動草稿位置；按 Esc 退出工具，確認後續單擊只用於選擇。

\subsection{測量週期結構中的區域性幾何}\label{sec:periodic-measure-reference}

讀取數值時，使用三維檢視工具欄的\ui{Measurements}選單，選擇測量型別後按提示依次點選。頂部\ui{測量幾何}完成後顯示\ui{Geometry measurement completed.}，不在提示框中列出讀數。檢視器還提供晶胞引數、鍵統計、點到面距離、兩平面夾角、配位數和最近鄰。下表說明檢視器中的選取順序；可根據要測量的幾何量選擇工具。

\begin{table}[htbp]
\centering\small
\begin{tabularx}{\linewidth}{lX}
\toprule
任務 & 選取方式與讀數 \\
\midrule
兩點距離 & 按順序取兩個原子，距離用\AA。 \\
鍵角 & 取三個原子，第二個為頂點，角度用度。 \\
二面角 & 依次取四個原子，中間兩個確定轉軸，注意讀數符號。 \\
點到面距離 & 先取待測點，再取三個定義平面的點；平面內三點不可共線。 \\
兩平面夾角 & 依次取第一平面的三個點，再取第二平面的三個點。 \\
配位數或最近鄰 & 選中心原子，檢視當前鄰接條件下的數目或距離列表。 \\
晶胞引數 & 從選單直接讀取當前晶胞，無須選原子。 \\
鍵統計 & 選兩個原子以指定元素對，讀取當前鄰接演算法找到的該元素對的鍵數、平均和最短最長距離。 \\
\bottomrule
\end{tabularx}
\caption{按問題選擇幾何測量方式。}
\end{table}

頂部\ui{測量幾何}的返回資料在週期結構中使用鄰近映像連線區域性幾何；檢視器中的互動測量則按照實際點選的顯示映像連線，選擇哪一個邊界映像會影響讀數。最近鄰列表會列出位點與映像偏移，讀數時應一起檢視。配位數還取決於鄰接判定，不是對所有材料通用的固定化學常數。

要修改測得的幾何量，可在相應測量編輯中輸入目標值，也可使用頂部的\ui{設定距離}、\ui{設定鍵角}、\ui{設定二面角}。分子的\ui{atom}只動目標原子，\ui{subtree}與\ui{fragment}按連線尋找可移動部分；環不能被一根鍵分開時，移動範圍會受限。在週期晶體中，這些精確幾何命令只移動目標位點，不會自動把整段吸附分子當作分子子樹。若要整體平移吸附物，應選擇其全部位點，改用\ui{組分變換}。

\subsection{儲存和連線常用片段}\label{sec:fragment-library-reference}
\index{Pianduanku@片段庫}

片段草繪器把常見區域性結構組織為可檢索的庫。片段的連線埠說明它怎樣與宿主連線；有的片段適合單端連線，有的埠需要多個宿主位點，非共價埠則不提供這裡的共價連線按鈕。

\begin{steps}
  \item 在宿主分子中先選連線位點，再從\ui{分子建模}選單開啟\ui{片段草繪器}。搜尋框可以縮小列表；選中條目後檢視右側三維預覽、原子數、來源和標籤。
  \item 在\ui{連線模式}中選擇合適埠，閱讀連線說明，核對宿主的選擇數量。若\ui{連線片段}為灰色，先檢查是否選了不適用於共價連線的埠。
  \item 單擊\ui{連線片段}，檢查新增部分的位置與宿主連線。內建條目的連線頭已有定義，不宜僅憑預覽中的球大小猜測連線原子。
  \item 應用後核對宿主與片段的連線、氫數及連線處的距離，再執行\ui{分子診斷}並儲存專案。需要另一種連線方式時，撤銷後重新選埠，避免在錯誤連線上繼續疊加修改。
\end{steps}

要把自己的區域性結構放入庫，先讓當前分子文件只包含準備複用的完整片段，檢查連線和原子順序，再開啟\ui{儲存使用者片段}，填寫名稱。這個命令儲存當前分子，不是僅儲存螢幕中選亮的幾個球。要替換已有同名條目時明確勾選覆蓋，否則換一個名稱。

片段草繪器中的\ui{匯入片段庫}與\ui{匯出片段庫}用於交換片段庫 JSON。匯出庫檔案便於在另一臺電腦匯入，與匯出當前模型結構是兩件事。匯入後重新搜尋條目並檢查來源、埠和預覽；完成連線後的模型仍應單獨儲存。

\section{處理週期晶體}

\subsection{改變晶胞而保留所需幾何}\label{sec:lattice-transform-reference}

晶格矩陣的三行分別為 $\mathbf a,\mathbf b,\mathbf c$，三列為全域性 $x,y,z$ 分量，單位為\AA。\ui{編輯晶格}中的矩陣是可編輯表格，應逐單元格填寫；它不是把九個數一次輸入一個文字框。

設編輯前後的晶格分別為 $A$、$A'$。\ui{保持笛卡爾座標不變}決定了原子怎樣適應新外框，其關係可直接從式\eqref{eq:row-lattice-coordinates}讀出：
\begin{equation}
 \begin{aligned}
 \text{保留分數座標：}\quad&\mathbf f'=\mathbf f,
       &\mathbf r'&=\mathbf f A',\\
 \text{保留笛卡爾位置：}\quad&\mathbf r'=\mathbf r,
       &\mathbf f'&=\mathbf r(A')^{-1}.
 \end{aligned}
 \label{eq:lattice-preservation-options}
\end{equation}
例如，立方晶胞邊長由4改為8\,\AA，一個原子的原分數座標為 $(0.5,0,0)$，原位置為 $(2,0,0)\,\text{\AA}$。不勾選保持笛卡爾位置時，它仍在晶胞的一半處，位置變為 $(4,0,0)\,\text{\AA}$；勾選時，位置保留為 $(2,0,0)\,\text{\AA}$，分數座標改為 $(0.25,0,0)$。兩種選擇都可能有用，先決定要保留哪一種幾何關係，再改表格。

\begin{table}[!htb]
\centering\small
\begin{tabularx}{\linewidth}{p{30mm}XX}
\toprule
操作 & 保留什麼 & 重點檢查什麼\\
\midrule
編輯晶格，保留分數座標 & $\mathbf f$和位點數 & 胞長、夾角及改變後的原子間距。\\
編輯晶格，保留笛卡爾位置 & $\mathbf r$和位點數 & 新分數座標與週期邊界接觸。\\
構建超胞 & 原材料的週期排列 & $|\det S|$倍的體積和獨立位點數。\\
旋轉晶格 & 整體內部幾何、位點數 & 晶格與原子是否一起轉動。\\
交換座標軸 & 笛卡爾幾何、位點數 & 新基矢順序與對應分數座標。\\
映象晶格 & 分數座標、位點數 & 所選基矢取反及晶胞手性。\\
施加應變 & 分數座標、位點數 & 三條基矢的長度及受影響的距離。\\
\bottomrule
\end{tabularx}
\caption{按希望保留的量選擇晶格操作。表中的內部幾何指原子間距離與角度；改變週期邊界時，還需結合新晶胞檢查鄰近映像。}
\label{tab:round10-lattice-invariants}
\end{table}


\begin{steps}
  \item 儲存原始晶胞，開啟\ui{晶格編輯}中的\ui{編輯晶格}。先記下原胞長、夾角和代表性原子的笛卡爾位置。
  \item 逐格輸入新矩陣。要讓原子隨晶格形變，保留分數座標；要改變外框而保持當前笛卡爾位置，選擇\ui{保持笛卡爾座標不變}。例如，增大盒子裝入孤立結構時，通常更關心後一種幾何關係。
  \item 預覽檢查原子是否按預期移動，尤其檢查跨邊界的近鄰。應用後核對晶胞體積與位置，再另存。
\end{steps}

\ui{重定義晶胞}直接使用新晶格矩陣並保留原子的笛卡爾位置；它不按新舊晶胞體積比複製原子。需要真正增加原子數時，應使用\ui{構建超胞}。

\subsubsection{反向一條晶格向量}

\ui{映象晶格}輸入晶格向量序號1、2或3，將相應向量反向，其他兩行不變。用一個對角矩陣 $D_i$ 表示這一操作：它的第 $i$ 個對角元素為 $-1$，另外兩個為1，於是
\begin{equation}
 A'=D_iA,\qquad \mathbf f'=\mathbf f.
 \label{eq:mirror-lattice-row}
\end{equation}
例如選第1條軸，就是 $\mathbf a'=-\mathbf a$。分數座標 $(u,v,w)$ 對應的位置由 $u\mathbf a+v\mathbf b+w\mathbf c$ 變成 $-u\mathbf a+v\mathbf b+w\mathbf c$。由於 $\det D_i=-1$，晶胞體積的大小不變，基矢的手性改變。

這與\ui{映象原子}的作用不同。在一般斜晶胞中，反向 $\mathbf a$ 並不同時把 $\mathbf b$、$\mathbf c$ 關於某個笛卡爾平面反射，原子間的某些距離可能改變。若目的是將整個分子關於空間中的平面翻到另一側，應使用前面的平面映象公式；若要改晶格基矢方向，則按本式操作並核對原子位置。

\ui{旋轉晶格}輸入旋轉角度和全域性軸方向，會同時改變晶格與隨晶格擺放的原子；其行矩陣關係為 $A'=AR^{\mathsf T}$，分數座標保持不變。\ui{交換座標軸}則同時重排晶格行與分數座標分量：依次填2、1、3可交換 $\mathbf a$ 與 $\mathbf b$，同時交換 $u$ 與 $v$，因此笛卡爾位置不變。每次只做一種變化，按表\ref{tab:round10-lattice-invariants}核對後再儲存。

\subsubsection{沿晶格方向施加應變}

\ui{施加應變}和\ui{結構應變}使用無量綱相對伸縮量。為突出作用方向，本書將三個輸入記為 $\varepsilon_a,\varepsilon_b,\varepsilon_c$；它們依次縮放晶格矩陣的三行：
\begin{equation}
 A'=\begin{bmatrix}
 1+\varepsilon_a&0&0\\
 0&1+\varepsilon_b&0\\
 0&0&1+\varepsilon_c
 \end{bmatrix}A,\qquad \mathbf f'=\mathbf f.
 \label{eq:lattice-axial-strain}
\end{equation}
對話方塊中的三個分量仍以 $x,y,z$ 命名；在這個工具中應按第1、2、3條晶格向量理解，斜胞時不能把它們當作沿全域性 $x,y,z$ 分別拉伸。標量0.01使三行都乘1.01；輸入\literalinput{0.01 0 0}只使 $\mathbf a$ 乘1.01。0.01表示伸長1\%，不是增加 $0.01\,\text{\AA}$，也不是將胞長直接設為1.01\,\AA。

仍用第\ref{ch:modeling}章的手算斜胞，$\mathbf a=(4,0,0)$、$\mathbf b=(1,3,0)$、$\mathbf c=(0,0,6)$，單位為\AA。施加\literalinput{0.01 0 0}後，只有 $\mathbf a$ 變為 $(4.04,0,0)\,\text{\AA}$；$\mathbf b$ 的 $x$ 分量仍為1\,\AA。分數座標 $(0.5,0.5,0)$ 不變，笛卡爾位置則由 $(2.5,1.5,0)$ 變為 $(2.52,1.5,0)\,\text{\AA}$。這個結果也說明，它不是把所有原子的全域性 $x$ 座標統一乘1.01。

先儲存未變形模型，再填寫應變、預覽並應用。按上述小比例檢查基矢長度和代表性原子的位置，最後另存結果。這兩個工具提供三條晶格向量的伸縮，不含剪下應變；施加形變後若需要平衡結構，應另行設定相應的最佳化計算。


\subsection{建立正交超胞}\label{sec:orthogonal-cell-reference}

有些後續操作希望晶胞角接近 $90^\circ$。\ui{晶胞正交化}搜尋與原晶格相容的整數超胞，不是任意把斜晶胞拉成直角。模型可能變大，原子數和體積應一起檢查。

\begin{steps}
  \item 從原始小晶胞開始，在\ui{超胞與晶格}中開啟\ui{晶胞正交化}。
  \item 填寫最小、最大胞長，單位為\AA；最大值必須大於最小值。再給定原子數上限和角度容差，容差以度表示。介面會據當前晶格提供起始長度範圍，預設原子數上限5000、直角容差 $0.1^\circ$。
  \item 先用適合自己計算規模的原子上限預覽。完成後檢查三個角與 $90^\circ$ 的差、實際胞長以及位點數，而不是隻看外框在當前相機下像不像矩形。
  \item 若無候選，依次擴大最大胞長、適當增加原子上限或放寬角度容差，保留每次使用的條件。取消失敗或不合適的預覽，回到原模型。
  \item 應用後儲存為有區別的檔名，並記下正交化後的獨立位點數。計算輸入應引用這個新模型。
\end{steps}

\subsection{檢查對稱性並選擇標準晶胞}\label{sec:symmetry-reference}

\ui{分析對稱性}根據位置和晶格判斷空間群。它不修改結構；\ui{原胞}與\ui{慣用胞}則會生成標準表達，位點數和座標可能改變。名稱相近的\ui{轉為 P1}用於去除對稱性約束，適合開始區域性編輯，不能用它代替原胞約化。

\begin{steps}
  \item 在\ui{對稱性工具}中選擇\ui{分析對稱性}，從位置容差 $0.01\,\text{\AA}$、角度容差 $5^\circ$ 開始檢視。讀取空間群符號、Hall符號和國際編號，儲存到自己的建模記錄。
  \item 對有輕微數值偏差的輸入，可逐次比較更細或更粗的位置容差。容差用於判斷近似等價；較大的容差可能合併真實畸變，應結合結構來源選擇。
  \item 要比較原胞與慣用胞，從同一原始模型分別執行兩項構建，各自另存。比較化學式、體積、位點數及元素比例，確認它們表達相同的週期材料。
  \item 要對晶格本身作空間觀察，開啟\ui{Wigner-Seitz 晶胞}。結果顯示由晶格定義的多面體；它不把原子裁成一個團簇，也不改變當前模型位點。需要結構檔案時仍儲存原晶體。
\end{steps}

開始區域性缺陷編輯前，可再練習一次 P1 轉換。開啟按空間群構建後保留的\file{Cu-fcc}模型，在\ui{對稱性工具}中選擇\ui{轉為 P1}；這項命令沒有額外數值引數。執行後核對四個 Cu 位點和座標仍與原模型一致，再另存為\file{Cu-P1}。如果隨後只改一個區域性位置，應在這個副本上繼續，保留原有空間群模型供比較。P1 表示不再施加額外的空間群約束，不會自動讓原來規則的排列變得不規則。

\subsection{明確哪些原子允許移動}\label{sec:fixed-sites-reference}

\ui{固定原子}為模型寫入選擇性移動標記。它用於說明後續結構弛豫的約束，不能僅靠三維檢視中某個原子看起來不動來判斷約束是否存在。

\begin{steps}
  \item 選中準備固定的全部位點，開啟\ui{原子編輯}中的\ui{固定原子}，核對原子序號。
  \item 若薄片帶有幾何層資訊，介面還會出現固定底層數。要固定最低若干層時輸入層數；要嚴格使用自己的位點列表，將底層數設為0。
  \item 應用後儲存專案，並匯出原生結構JSON。重新開啟該結構，核對位點數與準備固定的序號。若需要直接檢查記錄，可用文字編輯器只讀檢視這些位點的\file{selective_dynamics}：三個false表示三個方向都固定，三個true表示都可移動。
\end{steps}

該命令將所選位點設為三個方向均固定，其餘位點設為可移動。若需同時固定兩組原子，先把兩組合成完整選擇再執行。儲存選擇集可以減少重複選取時的遺漏。

在分子建模的\ui{幾何最佳化}視窗，約束由\ui{固定原子索引}填寫；例如輸入\literalinput{1 2}可固定兩個指定原子，再預覽最佳化結果。本節的晶體固定標記用於記錄位點約束，不能用它替代計算任務本身的約束設定。第~\ref{ch:workflows}~章的氫分子最佳化不含固定原子；需要帶約束計算時，先按所用計算方法確認約束輸入，再檢查最佳化前後固定原子的座標。

\subsection{建立空位以外的點缺陷}\label{sec:other-defects-reference}

\ui{缺陷與合金}\,$\rightarrow$\,\ui{建立點缺陷}還包含替位、間隙和反位。它們改變原子的方式不同，應從預先儲存的宿主副本分別建立，避免把幾個缺陷無意疊在一起。

\begin{steps}
  \item 建立替位缺陷時，選擇\ui{替位}，指定一個原子序號與新元素。應用後位點數應保持不變，組成中相應元素的數目發生一增一減。
  \item 建立間隙缺陷時，選擇\ui{間隙}，輸入元素，明確選擇分數或笛卡爾座標，再填位置。預覽檢查最近鄰接觸，應用後位點數應增加一個。間隙位置應結合晶格空隙與最近鄰距離確定。
  \item 建立反位缺陷時，選擇\ui{反位}，填寫兩個位點序號。這會交換二者的物種，保留位置與總組成；兩個相同元素互換不會產生有意義的化學反位。
  \item 每個結果單獨命名，記錄宿主、超胞大小與缺陷位置。改變原胞表達或排序後，重新核對序號。
\end{steps}

\ui{列舉點缺陷}用於空位或替位候選，輸入允許考慮的位點範圍、替換元素、參考電荷和對稱容差。結果表中的\ui{Degeneracy}表示等價候選數，\ui{NetCharge}與輸入的參考電荷有關，最短接觸用於幾何檢查；它們不等於缺陷形成能或已經求得的電子電荷態。應用列舉後，各缺陷結構自動加入專案瀏覽器的結構分組，並開啟相應文件；結果表用於檢視，不在表內切換模型。到專案瀏覽器逐個開啟候選，核對組成與位點數，再儲存專案。後續比較採用一致的宿主與數值條件。

\subsection{生成顯式佔位的合金模型}\label{sec:sqs-reference}

SQS（特殊準隨機結構）在有限晶胞中安排具體的元素位置，使鄰近元素配對等統計量儘量接近隨機混合。它把輸入的混合佔位轉成每個位點由一種元素完整佔據的模型，便於繼續計算；這不表示材料已經變為化學有序相。命令讀取的是輸入結構中的元素與佔位比例，不會根據純元素晶體自動決定合金成分。

下面的四位點銅鎳例項給出輸入檔案、倍率、搜尋時限和結果核對。換用其他組成時，先選擇能容納整數原子個數的晶胞；例如等比例兩組分需要能分成相等的整數數量。儲存時同時保留混合佔位輸入和生成後的具體排列。

\subsection{示例：建立四位點銅鎳\headinglatin{SQS}}\label{sec:sqs-cu-ni-example}
\index{SQS@SQS 模型}

隨書\file{sqs-cu-ni-input.json}提供一個可以直接開始的輸入：邊長為 $3.61\,\text{\AA}$ 的面心立方晶胞，含4個獨立位點。每個位點均記錄Cu佔位0.5、Ni佔位0.5，總組成相當於Cu$_2$Ni$_2$。這裡的混合佔位描述統計上的組分分佈，並不是在同一位置放置兩個實際原子。SQS將這些比例轉成明確的原子種類，使每個位點由一種元素完整佔據。

本例採用示例晶格常數和較小晶胞，便於觀察混合佔位如何轉為顯式排列。研究合金性質時，可進一步擴大模型，並檢查關聯函式與晶格引數的收斂。

\begin{table}[htbp]
\centering\small
\begin{tabularx}{\linewidth}{lX}
\toprule
專案 & 本例設定與檢查 \\
\midrule
輸入檔案 & \file{sqs-cu-ni-input.json}；另附同名CIF。\\
晶格 & $a=b=c=3.61\,\text{\AA}$，三個夾角均為 $90^\circ$。\\
獨立位點 & 4；分數座標為 $(0,0,0)$、$(0,\frac12,\frac12)$、$(\frac12,0,\frac12)$、$(\frac12,\frac12,0)$。\\
位點佔位 & 每個位點Cu 0.5、Ni 0.5；兩種佔位之和為1。\\
SQS 超胞倍率 & 在文字框中輸入\literalinput{1}，保持本例四位點晶胞。\\
SQS 搜尋時限 & 輸入1秒；本例進入小規模精確列舉。\\
預期結果 & Cu 2個、Ni 2個，共4個位點，每個位點佔位為1。\\
\bottomrule
\end{tabularx}
\caption{四位點銅鎳SQS練習。}
\end{table}

\begin{steps}
  \item 將\file{sqs-cu-ni-input.json}匯入專案，開啟其三維文件。線上版中先上傳檔案，再從雲端位置匯入。將文件命名為\file{CuNi-mixed}，儲存起始專案，以便隨時回到混合佔位輸入。
  \item 檢視獨立位點數、晶胞和原子資訊。位點數應為4，位點組分應同時包含Cu和Ni，各為0.5。若檢視輸入JSON，每個\file{sites}條目中的\file{species}都有兩個元素及各自的\file{occu}值；這比僅憑顯示顏色判斷組成更可靠。
  \item 保持該晶體文件啟用，從\ui{晶體建模}\,$\rightarrow$\,\ui{缺陷與合金}開啟\ui{生成 SQS 模型}。在\ui{SQS 超胞倍率}中輸入\literalinput{1}，在\ui{SQS 搜尋時限（秒）}中輸入1。
  \item 單擊\ui{預覽}，檢視四個位點是否已分配成兩種明確的元素，確認晶胞大小和位點數沒有增加，再\ui{應用}。本例保留原有晶格和位置，改變的是各個位點的物種佔據。
  \item 檢查結果組成：應為Cu$_2$Ni$_2$，位點數仍為4，每個位點只含一種元素、佔位為1。測量最近鄰距離，應約為 $3.61/\sqrt{2}=2.552655\,\text{\AA}$。元素分配可以透過晶胞對稱操作得到等價表達，檢查時以組成、佔位和週期幾何為準。
  \item 按前述結構匯出方法匯出原生JSON，命名為\file{sqs-cu-ni-ordered.json}，保留搜尋記錄。檔名中的\file{ordered}表示每個位點已有明確元素，不表示生成了化學長程有序相。
  \item 將結果文件命名為\file{CuNi-SQS-4}並儲存專案；再匯出CIF，便於交換晶胞和明確的元素佔位。線上版中下載匯出檔案，並保留輸入JSON、結果JSON和專案副本。CIF用於結構交換，SQS搜尋記錄以原生JSON為準。
\end{steps}

\subsubsection*{檢視詳細搜尋記錄}

用文字編輯器開啟匯出的JSON，只檢視\file{properties.sqs}中的記錄，無需修改檔案。\file{search_mode}表示搜尋方式，\file{configuration_count}表示候選排列數，\file{objective_function}是關聯匹配的目標值。

參考結果中的\file{search_mode}為\literalinput{exact-enumeration}，\file{configuration_count}為6，\file{objective_function}約為0.1666667。六種排列來自四個位點中選擇兩個放Cu，其餘放Ni。本例六對不同位點的週期最短距離相同；Cu--Cu與Ni--Ni各佔其中一對，有限晶胞不能完全再現獨立隨機佔位的配對比例，因此目標值並不為零。這個無量綱目標值衡量當前關聯函式與隨機佔位目標的差別，不是吸附能、總能量或缺陷形成能。

搜尋時限控制候選搜尋階段；完整操作還包括精確列舉、等價結構比較及結果整理，因此總耗時可能超過設定值。擴大模型時，同時考慮位點數與候選規模。

\Needspace{10\baselineskip}
\begin{notebox}{倍率與混合佔位}
本命令的單個整數倍率表示晶胞體積倍數，不是三個方向各乘同一個數。要明確三個方向的重複數，在同一倍率文字框中輸入例如\literalinput{2 2 2}；它會把本例擴為32個位點。比較模型前應記錄實際晶胞、組成和搜尋條件。混合佔位輸入用於建模；電子結構計算應使用生成並檢查過的顯式佔位結果。匯入混合佔位時，計算物件可能給出分數佔位不能直接表示的提示，不能把轉換出的單一元素顯示當成SQS已經完成。
\end{notebox}

隨書\file{sqs-cu-ni-ordered.json}和\file{sqs-cu-ni-ordered.cif}提供對應結果；\file{sqs_example.m}可在以\file{-nojvm}啟動的MATLAB\regmark{}工具箱環境中重新生成輸入和結果。生成顯式佔位模型後，可繼續建立電子結構計算。


\subsection{選擇低維結構的構建方式}\label{sec:nano-reference}

第~\ref{sec:nanostructure-examples}~節已用石墨烯和銅完成五類構建。這裡按希望保留的邊界選擇入口，各例引數見該節。

\begin{table}[htbp]
\centering
\begin{tabularx}{\linewidth}{lX}
\toprule
需要的模型 & 命令與關鍵區別 \\
\midrule
有限寬度的片層 & \ui{擷取奈米帶}；按寬度、長度重複數裁切，沿帶長週期。\\
捲曲的片層 & \ui{生成奈米管}；按手性指數與軸向重複數構建，沿管軸週期。\\
有限橫截面的晶體 & \ui{擷取奈米線}；按晶格方向與半徑擷取，保留軸向週期。\\
三個方向均有限的團簇 & \ui{量子點 / 奈米空洞}的\ui{量子點}模式；保留球內原子。\\
週期塊體中的球形孔洞 & 同一視窗的\ui{奈米空洞}模式；去掉球內原子，保留外圍週期宿主。\\
\bottomrule
\end{tabularx}
\caption{按邊界選擇奈米結構命令。}
\end{table}

線、帶和管的重複數不是以\AA 為單位的總尺寸。比較不同模型時，以實際晶胞、位點數和週期方向為準；需要鈍化時，先儲存未鈍化輸入，再檢查新增原子的組成與邊界配位。

\FloatBarrier
\section{建立軟物質模型}

\subsection{構建帶有支鏈的高分子}\label{sec:branched-polymer-reference}

支化構建從一條主鏈出發，按重複單元間隔連線片段，適合學習主鏈與支鏈的拓撲關係。重複單元名稱來自重複單元庫，支鏈名稱來自片段庫，兩者應分別核對。

\begin{steps}
  \item 先匯入隨書\file{user-repeat-ethane-input.json}並啟用分子文件，從\ui{分子建模 → 軟物質 → 聚合物建模}開啟\ui{構建支化聚合物}。保留輸入副本；構建結果會替換活動文件中的分子。先用短鏈練習，例如主鏈重複單元\literalinput{PE}、重複數8、支鏈\literalinput{Methyl}、支化間隔3、種子1。
  \item 檢視預計原子規模，給出適合練習的原子上限。支化間隔按重複單元計數，不是按笛卡爾長度或主鏈上每一個原子計數。
  \item 預覽並檢查支鏈是否連在預期主鏈單元附近，再確認主鏈的兩端與連線。若需改變支鏈密度，回到原始短鏈調整間隔，不要在已經支化的結果上反覆執行構建。
  \item 應用後檢查組成、支鏈連線和區域性距離，再執行\ui{分子診斷}並儲存原生模型。需要建立指定密度的材料時，再進行分子裝填。
\end{steps}

\ui{構建均聚物}可用於不含支鏈的對照：選一個重複單元、重複數、獨立鏈數、端基與初始構象，固定種子後構建。獨立鏈數大於1表示生成多條鏈，不是自動建立分子間共價鍵。\ui{伸展}、\ui{折線}與\ui{無規線團}是初始幾何選擇；儲存後仍應檢查鏈內接觸及後續計算是否需要進一步整理。

\subsection{建立基礎樹枝狀骨架}\label{sec:dendrimer-reference}

\ui{構建基礎樹枝狀聚合物}按代數與分支數展開骨架，適合觀察拓撲層次。核心官能度決定第一代分支數，分支因子決定後續每個分支繼續產生多少分支。骨架元素和鍵長定義初始幾何，實際材料的官能團與化學組成可在此基礎上繼續完善。

\begin{steps}
  \item 匯入隨書\file{user-repeat-ethane-input.json}並啟用，保留輸入副本。從\ui{分子建模 → 軟物質 → 聚合物建模}開啟\ui{構建基礎樹枝狀聚合物}，先用代數2、核心官能度3、分支因子2、元素C、骨架鍵長 $1.54\,\text{\AA}$、種子1作小模型練習。
  \item 檢視預計規模，設定原子數上限。增加代數會使分支快速增多，先看小模型的連線再擴充套件。
  \item 預覽時找到中心，順著連線檢視第一代和第二代。檢查分支是否意外穿過其他分支，檢視最短接觸。
  \item 應用後檢查組成、分支連線和區域性接觸，執行\ui{分子診斷}後儲存專案和原生結構。得到初始骨架後，可繼續整理幾何並進行弛豫。
\end{steps}

\subsection{儲存自定義重複單元}\label{sec:repeat-unit-reference}

重複單元除了原子和鍵，還需要兩個不同的連線端。頭原子用於連線前一個單元，尾原子用於連線後一個單元；離去原子在連線前移除。下面用乙烷中的兩個碳定義一個簡單骨架，再連線成二聚體，檢視這些欄位怎樣影響實際結構。

隨書\file{user-repeat-ethane-input.json}含兩個C和六個H，共8個原子、7條單鍵。C為1、2號，H為3至8號。輸入中的C--C長度為$1.52\,\text{\AA}$。本例移除六個H後儲存碳骨架，構建鏈時再由工具統一補氫；這是建立單元的幾何操作，不用於表示乙烷的聚合反應。

\begin{table}[htbp]
\centering
\begin{tabularx}{\linewidth}{lX}
\toprule
操作欄位 & 本例填寫 \\
\midrule
使用者重複單元名稱 & \literalinput{MANUAL_CC}。\\
頭原子、尾原子 & 分別為1、2。\\
離去原子序號 & 在文字框輸入\literalinput{3 4 5 6 7 8}。\\
均聚物重複單元 & 輸入儲存的名稱\literalinput{MANUAL_CC}。\\
重複次數、鏈數 & 分別設為2、1。\\
初始構象、種子 & 選擇\ui{伸展}，種子7。\\
首尾端基 & 均保持\literalinput{none}；其餘未滿足價鍵由氫補足。\\
原子數上限 & 1000；足以容納本例模型。\\
\bottomrule
\end{tabularx}
\caption{從自定義碳骨架建立二聚體。}
\end{table}

\begin{steps}
  \item 匯入\file{user-repeat-ethane-input.json}，開啟分子文件。逐個檢視前兩個原子，確認1、2號是碳；再核對3至8號均為氫。保留此輸入檔案，不在建立重複單元前重新排序。
  \item 從\ui{分子建模 → 軟物質 → 聚合物建模}開啟\ui{儲存使用者重複單元}，按表輸入名稱、頭尾及離去原子序號。頭尾必須是不同的現有原子，且不能同時列為離去原子。
  \item 首次儲存不勾選覆蓋。若已有同名示例條目，確認是本例舊記錄後再覆蓋，或改用另一個名稱並在後續保持一致。儲存後當前乙烷仍是8個原子；儲存單元本身不修改來源分子。
  \item 仍在這個分子文件中開啟\ui{構建均聚物}。填入儲存的名稱，按表設定重複次數、鏈數、構象、端基和種子，頭尾連線方式保持預設。預覽後應用，檢查生成的完整鏈。
  \item 結果應有4個C、10個H，共14個原子、13條單鍵，兩個單元之間新增一條C--C連線。旋轉檢查主鏈與末端氫，對可疑接觸測量距離，再執行\ui{分子診斷}。隨書\file{user-repeat-CC-dimer.json}提供對應的二聚體。
  \item 將結果命名為\file{MANUAL-CC-dimer}，儲存專案並匯出原生結構JSON。儲存輸入乙烷、重複單元引數和二聚體三個階段；只交換二聚體檔案不會把自定義單元自動註冊到另一臺電腦的庫中。
\end{steps}

用文字編輯器開啟輸入或結果JSON，可以檢視元素、座標和連線。隨書\file{user_repeat_unit_example.m}重複執行相同的儲存與構建，並將示例庫單獨寫到\file{user-repeat-unit-library.json}，不覆蓋個人重複單元庫。配套指令碼還提供組成、連線及診斷結果的檢查方法。繼續建立更長鏈時，從相同單元開始增加重複次數，再檢查完整鏈內的接觸。

\subsection{將一種分子裝入受限區域}\label{sec:confined-packing-reference}

\ui{構建無定形晶胞}以當前分子為組分，按分子數和密度放入週期盒；\ui{構建受限分子層}進一步限定可用區域。密度單位為 $\mathrm{g/cm^3}$，最小接觸用\AA。裝填得到的是初始位置，後續物理平衡需要單獨完成。

\begin{steps}
  \item 開啟連線與幾何已檢查的分子。先儲存組分原件，再從\ui{無定形建模}選擇單組分裝填或受限分子層。
  \item 輸入分子數、目標密度、最小接觸、種子與原子上限。首次練習使用較少分子和較低初始密度，先確認組分及邊界的處理方式。
  \item 對受限分子層，選擇限制軸1、2或3，再填寫區間上下限。例如 \literalinput{0.2 0.8} 表示該軸盒長的20\%至80\%；當兩個數都在0與1之間時按比例解釋，不能用它表達 $0.2$ 到 $0.8\,\text{\AA}$。
  \item 預覽時從限制軸方向的側面檢查分子是否落在區域內，確認沒有把層的空白部分誤當作可裝填體積。
  \item 應用後核對分子數、總原子數、盒長以及跨週期接觸，另存結果與輸入條件。
\end{steps}

\ui{密度遞增階段（可選）}可填寫一串正的、非遞減的密度，最後一個值必須等於目標密度；留空則只用目標密度。受限層練習中，密度遞增階段留空，以免與區域限制衝突。\ui{每批裝填分子數}控制裝填批次，不改變總分子數。需要按兩種分子組成裝填時，使用前面的\ui{裝填混合物}教程；摩爾分數和質量分數模式還需給出總分子數，最終應檢查換算後的整數數量。

\subsection{向已有盒子裝填分子}\label{sec:host-packing-reference}

已有宿主時，\ui{裝入已有晶胞}保留當前晶胞與宿主原子，把另一個專案分子加入盒中。它不根據目標密度重新建立一個立方盒。\ui{繞奈米粒子裝填}在此基礎上，為當前宿主周圍留出指定間隙。

\begin{steps}
  \item 把宿主和待裝填分子匯入同一專案，分別命名。啟用具有正交晶胞的宿主文件，在\ui{宿主裝填}中選擇需要的命令；斜晶胞需先整理為適合裝填的盒子。
  \item 在\ui{待裝填分子}中選擇組分，填寫分子數、最小接觸、種子、每批分子數和原子上限。這裡沒有透過密度重新計算盒長的步驟，應先保證宿主盒中有足夠空間。
  \item 繞奈米粒子裝填時再填寫粒子間隙，單位為\AA。排除區域以宿主平均位置為中心、覆蓋宿主最外原子的球形範圍再加間隙；它不逐個追隨任意凹凸表面。
  \item 預覽並從不同方向檢查宿主保留、分子位置與跨週期邊界的接觸。空間不足時，回到原宿主擴大適當的空白盒子，或減少裝填數量後重新開始。
  \item 完成後核對新增分子數與組成，將宿主、組分原件和裝填結果分別保留；儲存專案及原生結構。
\end{steps}

\section{完善表面模型}

\subsection{鈍化表面}\label{sec:surface-finishing-reference}

第~\ref{ch:advanced}~章的吸附練習介紹了候選位點與O--H放置。這裡處理薄片上仍需封端的邊界；從已儲存的薄片副本開始，使鈍化和後續溶劑處理分別留下記錄。

\begin{steps}
  \item 若表面有未飽和配位，需要新增封端原子，開啟\ui{表面鈍化}，選擇元素、目標鍵長、處理上表面/下表面/兩面，以及表面探測深度。
  \item 目標配位數為0時由工具推斷；已知配位目標時明確輸入。最小接觸用來排除過近位置。預覽中檢查兩側新增是否符合目的，並檢視原有原子有無被錯誤識別為表面。
  \item 應用後比較新增元素數與邊界配位，單獨儲存鈍化結果，再決定是否繼續裝填溶劑。
\end{steps}

\subsection{在薄片上方加入溶劑}

\ui{新增溶劑層}需要專案中已有一個完整溶劑分子。保持薄片文件啟用，在命令對話方塊的\ui{溶劑分子/結構}下拉選單中選擇它，再設定分子數、層厚、表面間隔、剩餘真空、最小接觸、種子和超時秒數。層厚、間隔、真空和接觸都用\AA；分子數為整數。先用少量分子檢視它們位於薄片哪一側，核對溶劑內部分子幾何、層內接觸和剩餘真空後再擴大。生成初始溶劑層後，再按所用計算方法設定溫度與平衡過程。

\section{複用建模過程}

\subsection{儲存常用的建模引數}\label{sec:modeling-preset-reference}
\index{Canshuyushe@引數預設}

引數預設儲存某個命令中填寫的條件。再次載入時，可根據材料和模型調整數值；生成的模型另行儲存。

\begin{steps}
  \item 開啟常用建模命令，填寫並核對引數。可以先預覽檢查效果；要儲存一組待使用的條件，也可在應用前儲存。
  \item 單擊對話方塊頂部的\ui{儲存預設}，輸入說明用途的名稱，例如“表面練習用厚度”。儲存後確認該名稱出現在\ui{預設}列表中。
  \item 關閉並重新開啟同一命令，選中名稱，再單擊\ui{載入}。逐項檢視引數是否符合當前模型，特別是位點序號、專案結構引用和長度。
  \item 必要時改動引數，再預覽、應用並儲存結果模型。載入預設本身不代表建模已經提交。
\end{steps}

對話方塊底部\ui{重置}恢復這個對話方塊的初始填寫值，同時清除臨時預覽；頂部歷史區的\ui{重置}則恢復當前結構草稿的基準，兩者作用不同。對話方塊中的\ui{幫助資訊}可展開檢視引數解釋，再收起繼續填寫。

\subsection{錄製可追溯的建模過程}\label{sec:modeling-recipe-reference}
\index{Jianmopeifang@建模配方}

Recipe記錄一系列已經提交的建模命令、引數、隨機種子和前後結構的校驗資訊。它適合保留“這個結果怎樣得到”的過程；原始模型與最終結構仍應同時儲存。

\begin{steps}
  \item 開啟並儲存一份明確命名的起始模型。啟用其三維文件，在\ui{歷史}中單擊\ui{錄製}。
  \item 執行需要記錄的建模操作。先預覽、檢查，再應用；取消的操作不寫入Recipe。
  \item 完成操作後，單擊\ui{儲存 Recipe}。錄製隨之停止，軟體按時間生成名稱，並在狀態資訊中給出儲存位置。
  \item 開啟\ui{歷史}\,$\rightarrow$\,\ui{預設與模板}，單擊\ui{重新整理}，在型別為\ui{操作配方}（Recipe）的條目中查詢本次記錄。儲存模型結果和專案，並把起始輸入一同歸檔。
\end{steps}

Recipe記錄當前結構文件中受支援的建模操作，範圍包括已應用的命令與引數。檢視變化、選擇集管理和分析檢視不寫入Recipe；從分析結果視窗應用候選後，另存最終模型。\ui{預設與模板}視窗提供條目檢視和模板建立。需要重放Recipe時，MATLAB\regmark{}工具箱使用者可採用建模API，並使用與記錄匹配的起始結構；圖形介面中可將Recipe保留為操作記錄。

\subsection{建立建模專案模板}\label{sec:modeling-template-reference}

\ui{預設與模板}把命令預設、Recipe和建模專案模板放在同一列表中，列出型別、名稱與詳情。這裡的建模專案模板記錄所用Recipe名稱及輸入輸出格式，與工作流頁面儲存的 \file{.wf} 計算模板用途不同。

\begin{steps}
  \item 在\ui{歷史}中開啟\ui{預設與模板}，先單擊\ui{重新整理}，確認準備引用的Recipe存在。
  \item 單擊\ui{新建專案模板}，填寫名稱、Recipe名稱、輸入格式和輸出格式。預設\literalinput{auto}表示由輸入識別格式，\literalinput{same}表示輸出沿用輸入格式；它們是格式選擇，不是檔案路徑。
  \item 儲存後在列表核對型別、名稱和詳情中的輸入到輸出關係。若條目標為無效，按詳情檢查相應記錄，確認所引用的記錄完整後再使用。
  \item 將模板和它引用的Recipe一起保留。實際模型檔案另行儲存；建立模板不會自動開啟一份新工程或執行建模。
\end{steps}

\subsection{檢視已有的批處理任務}\label{sec:modeling-jobs-reference}

\ui{建模任務}檢視已經儲存的建模檔案批處理任務，與電子結構工作流的SCF執行列表不同。視窗顯示編號、狀態、已完成數、請求數、成功數、是否可恢復以及錯誤資訊；一個批次完成，也可能存在個別檔案失敗。

\begin{steps}
  \item 開啟\ui{歷史}\,$\rightarrow$\,\ui{建模任務}，單擊\ui{重新整理}，按任務編號找到需要檢視的批次。
  \item 比較請求數、已完成數與成功數，再讀錯誤欄。其中成功數表示正常生成輸出的檔案數。
  \item 選中可恢復的條目，再單擊\ui{恢復任務}，等待狀態更新並重新整理列表。若按鈕為灰色，先確認該任務當前是否標為可恢復。
  \item 不準備繼續的可恢復任務可單擊\ui{取消任務}請求停止，隨後重新整理確認狀態。取消不是刪除已經寫出的結果，輸出檔案仍需按任務記錄檢查。
  \item 到任務原有輸出目錄檢查結果結構與報告，保留失敗條目的輸入，儲存輸入、輸出與任務記錄。
\end{steps}

這個視窗用於檢視與恢復已有任務。需要建立檔案批處理時，MATLAB\regmark{}工具箱使用者可使用建模API；圖形介面使用者按上述方法檢查已有任務。

\subsection{恢復尚未儲存的結構編輯}\label{sec:modeling-draft-reference}

撤銷與重做針對當前結構草稿；歷史區的\ui{重置}恢復草稿基準，儲存後基準隨模型更新。要長期比較兩個模型，應分別儲存命名結果。

重新開啟文件時若提示恢復草稿，核對對應結構後再恢復，並檢查原子數、組成、晶胞和最後編輯位置。另存結果；恢復不會重建中斷前的全部撤銷步驟。若不恢復，先確認原儲存檔案就是需要的版本，再清理草稿。
